\chapter{Blood Disorders}

\section*{Definition and Overview}
Blood disorders (also called haematological disorders) are conditions that affect the blood's main components: red blood cells, white blood cells, platelets, or the plasma proteins involved in clotting. The major categories are anaemia (too few red blood cells or too little haemoglobin), bleeding and clotting disorders (problems with platelets or clotting factors), and haematological cancers such as leukaemia, lymphoma, and myeloma. Blood disorders may be inherited or acquired and range from mild and easily treatable to serious and lifelong. This chapter provides a general overview; each specific disorder is covered in more detail in its own chapter.

\section*{Epidemiology}
Blood disorders are common worldwide. Anaemia affects roughly a quarter of the world's population, and iron deficiency is its most frequent cause, particularly in women of reproductive age and young children. Von Willebrand disease affects roughly one in 100 to one in 1,000 people, while the more severe haemophilia A affects about one in 5,000 to one in 10,000 males. Blood cancers are less common but significant, with leukaemia and lymphoma accounting for a substantial share of cancers diagnosed in children and younger adults. Many bleeding disorders are inherited and present in childhood, whereas acquired conditions such as the anaemia of chronic disease become more common with age.

\section*{Aetiology and Pathophysiology}
Blood cells are produced in the bone marrow from stem cells, and clotting depends on platelets acting with a cascade of plasma proteins. Disorders arise when any part of this system fails:

\begin{itemize}
    \item \textbf{Red cell disorders:} reduced production, increased destruction (haemolysis), or blood loss cause anaemia; inherited changes in haemoglobin cause sickle cell disease and thalassaemia
    \item \textbf{White cell disorders:} overproduction of abnormal white cells occurs in leukaemia, while underproduction from marrow failure or chemotherapy raises the risk of infection
    \item \textbf{Platelet and clotting disorders:} too few platelets or missing or abnormal clotting factors cause bleeding; an excess tendency to clot (thrombophilia) causes unwanted blood clots
    \item \textbf{Inherited causes:} gene changes passed through families, as in haemophilia, sickle cell disease, and thalassaemia
    \item \textbf{Acquired causes:} infections, nutritional deficiencies (iron, B12, folate), autoimmune disease, cancer, medicines, and treatments such as chemotherapy or radiotherapy
\end{itemize}

\section*{Clinical Features}
Symptoms depend on which component of the blood is affected, but common features include:

\begin{itemize}
    \item Fatigue, weakness, pallor, and breathlessness on exertion (anaemia)
    \item Easy bruising, nosebleeds, bleeding gums, or prolonged bleeding from cuts (platelet or clotting problems)
    \item Frequent or severe infections and slow healing (white cell problems)
    \item Fevers, night sweats, and unexplained weight loss (blood cancers)
    \item Swollen lymph nodes or an enlarged spleen
    \item Bone pain, especially in leukaemia and myeloma
    \item Jaundice or dark urine in haemolytic anaemia
\end{itemize}

\section*{Differential Diagnosis}
The same symptoms can be produced by many conditions, so a blood disorder must be distinguished from:

\begin{itemize}
    \item Chronic kidney disease, thyroid disease, and chronic infection or inflammation, which all cause anaemia
    \item Malnutrition and poor absorption, including coeliac disease and pernicious anaemia
    \item Solid cancers that spread to the bone marrow
    \item Liver disease, which impairs the production of clotting factors
    \item Drug-induced suppression of the bone marrow
    \item Other causes of bleeding, such as peptic ulcers or heavy menstrual bleeding
\end{itemize}

\section*{When to Seek Medical Care}
See a doctor for unexplained or persistent fatigue, easy bruising, unusual bleeding, fevers, night sweats, or weight loss, or if there is a family history of an inherited blood disorder.

\textbf{Call 000 (or local emergency services) for:}
\begin{itemize}
    \item Uncontrolled or heavy bleeding that does not stop with pressure
    \item Sudden shortness of breath or chest pain
    \item High fever with shaking chills in someone known to have a low white cell count
    \item Confusion, collapse, or any sudden, severe symptom
\end{itemize}

\section*{Diagnosis and Investigations}
The diagnosis of a blood disorder relies on blood tests interpreted together with the clinical picture:

\begin{itemize}
    \item \textbf{History and examination:} symptoms, family history, medicines, diet, and examination of the skin, lymph nodes, and spleen
    \item \textbf{Full blood count:} counts of red cells, white cells, and platelets
    \item \textbf{Blood film:} microscopic examination of the blood cells, looking at their size, shape, and maturity
    \item \textbf{Iron studies, B12, and folate:} to find the cause of anaemia
    \item \textbf{Coagulation tests:} to assess bleeding or clotting risk
    \item \textbf{Bone marrow biopsy:} when leukaemia, marrow failure, or myeloma is suspected
    \item \textbf{Genetic testing:} for inherited disorders and to guide the treatment of blood cancers
    \item \textbf{Imaging:} ultrasound, CT, or PET to assess lymph nodes and the spleen
\end{itemize}

\section*{Management}
Management is tailored to the specific disorder but draws on a limited set of strategies:

\begin{itemize}
    \item \textbf{Anaemia:} iron, vitamin B12, or folate replacement, together with treatment of the underlying cause
    \item \textbf{Bleeding disorders:} clotting-factor replacement (for example in haemophilia), desmopressin, or antifibrinolytic medicines
    \item \textbf{Blood transfusion:} to replace red cells, platelets, or plasma when counts are dangerously low
    \item \textbf{Blood cancers:} chemotherapy, targeted therapies, immunotherapy, and stem-cell (bone marrow) transplantation
    \item \textbf{Infection prevention:} vaccination and prompt treatment of infections, especially with low white cell counts
    \item \textbf{Supportive care:} physiotherapy, pain management, and psychological support, particularly for chronic or lifelong disorders
    \item \textbf{Regular follow-up:} monitoring of blood counts, treatment response, and complications
\end{itemize}

\section*{Complications}
\begin{itemize}
    \item Severe anaemia can strain the heart and worsen heart failure or angina
    \item Haemolysis can cause pigment gallstones and gall-bladder problems
    \item Bleeding into joints and muscles causes pain and progressive joint damage in severe haemophilia
    \item Blood cancers can spread or relapse and lead to serious infections and bleeding
    \item Repeated transfusions can cause iron overload, which damages the liver and heart
    \item Long-term treatments carry their own risks, including infertility and second cancers
\end{itemize}

\section*{Prognosis and Outcomes}
The outlook varies widely with the specific disorder. Most forms of anaemia are easily corrected once the cause is found. Inherited bleeding disorders cannot be cured but are usually well managed with replacement therapy, and most affected people live full lives. Outcomes in leukaemia, lymphoma, and myeloma have improved markedly with modern treatment, and many patients achieve long-term remission or cure. For chronic disorders, adherence to treatment and regular follow-up greatly improve quality of life and survival.

\section*{Prevention and Self-Care}
\begin{itemize}
    \item Eat an iron-rich, balanced diet and address heavy menstrual bleeding or other chronic blood loss
    \item Take iron, folate, or B12 supplements if a deficiency is diagnosed
    \item Avoid excess alcohol and medicines that harm the bone marrow
    \item Keep vaccinations up to date if you have a blood disorder or a low white cell count
    \item If you have a bleeding disorder, take precautions against injury and wear medical-alert identification
    \item Attend genetic counselling if there is a family history of an inherited disorder and you plan a family
    \item Have regular check-ups and follow your treatment plan
\end{itemize}

\section*{Sources and Further Reading}
The following sources provide reliable, up-to-date information for patients and students of medicine.

\begin{itemize}
    \item NHS. \textit{Overview of blood and bone marrow conditions.} https://www.nhs.uk/conditions/blood-bone-marrow/
    \item Mayo Clinic. \textit{Blood disorders -- symptoms, causes, and treatment.} https://www.mayoclinic.org/
    \item World Health Organization. \textit{Anaemia and nutritional deficiency -- fact sheets.} https://www.who.int/
    \item MSD Manual. \textit{Hematology and oncology -- professional edition.} https://www.msdmanuals.com/
    \item MedlinePlus. \textit{Blood disorders -- patient information.} https://medlineplus.gov/
    \item Centers for Disease Control and Prevention. \textit{Blood disorders information and resources.} https://www.cdc.gov/blood-disorders/
\end{itemize}
